\input{preamble}

\section*{Rutherford scattering}

Rutherford scattering is the interaction
$e^-+X\rightarrow e^-+X$ where $X$ is an atom.

\begin{center}
\begin{tikzpicture}
\draw[dashed] (0,0) circle (0.5cm);
\draw (0,0) node {$X$};
\draw[thick,->] (-2,0) node[anchor=east] {$e^-$} -- (-0.6,0);
\draw[thick,->] (0.40,0.40) -- (1.3,1.3) node[anchor=south west] {$e^-$};
\draw (1,0.5) node {$\theta$};
\end{tikzpicture}
\end{center}

These are the momentum vectors for Rutherford scattering where $E=\sqrt{p^2+m^2}$.
\begin{equation*}
p_1=\underset{\substack{\text{inbound}\\ \text{electron}}}
{\begin{pmatrix}E\\0\\0\\p\end{pmatrix}}
\qquad
p_2=\underset{\substack{\text{outbound}\\ \text{electron}}}
{\begin{pmatrix}
E\\
p\sin\theta\cos\phi\\
p\sin\theta\sin\phi\\
p\cos\theta
\end{pmatrix}}
\end{equation*}

Spinors for the inbound electron.
\begin{equation*}
u_{11}=\frac{1}{\sqrt{E+m}}
\underset{\substack{\text{inbound electron}\\ \text{spin up}}}
{\begin{pmatrix}E+m\\0\\p\\0\end{pmatrix}}
\qquad
u_{12}=\frac{1}{\sqrt{E+m}}
\underset{\substack{\text{inbound electron}\\ \text{spin down}}}
{\begin{pmatrix}0\\E+m\\0\\-p\end{pmatrix}}
\end{equation*}

Spinors for the outbound electron.
\begin{equation*}
u_{21}=\frac{1}{\sqrt{E+m}}
\underset{\substack{\text{outbound electron}\\ \text{spin up}}}
{\begin{pmatrix}E+m\\0\\p_{2z}\\p_{2x}+ip_{2y}\end{pmatrix}}
\qquad
u_{22}=\frac{1}{\sqrt{E+m}}
\underset{\substack{\text{outbound electron}\\ \text{spin down}}}
{\begin{pmatrix}0\\E+m\\p_{2x}-ip_{2y}\\-p_{2z}\end{pmatrix}}
\end{equation*}

The probability density $|\mathcal M_{ab}|^2$ for spin state $ab$ is
\begin{equation*}
|\mathcal{M}_{ab}|^2=\left(\frac{Z\alpha}{q^2}\right)^2\left|\bar{u}_{2b}\gamma^0 u_{1a}\right|^2
\end{equation*}

where $q^2$ is momentum transfer such that
\begin{equation*}
q^2=(p_1-p_2)^2=(p_1-p_2)^\mu g_{\mu\nu}(p_1-p_2)^\nu=2p^2(\cos\theta-1)
\end{equation*}

The expected probability density
$\langle\vert\mathcal{M}\vert^2\rangle$
is the average of spin states.
\begin{equation*}
\langle\vert\mathcal{M}\vert^2\rangle
=\frac{1}{2}\sum_{a=1}^2\sum_{b=1}^2\left|\mathcal{M}_{ab}\right|^2
\end{equation*}

The Casimir trick uses matrix arithmetic to sum over spin states.
\begin{equation*}
\langle\vert\mathcal{M}\vert^2\rangle
=\frac{1}{2}\left(\frac{Z\alpha}{q^2}\right)^2
\operatorname{Tr}\left[(\slashed{p}_1+m)\gamma^0(\slashed{p}_2+m)\gamma^0\right]
\end{equation*}

The result is
\begin{equation*}
\langle\vert\mathcal{M}\vert^2\rangle
=2\left(\frac{Z\alpha}{q^2}\right)^2\left(E^2+m^2+p^2\cos\theta\right)
\end{equation*}

Recall that $m\approx0.511\,\text{MeV}$.
For $m\gg p$ we can use the approximation
\begin{equation*}
E^2+m^2+p^2\cos\theta\approx2m^2
\end{equation*}

Hence
\begin{equation*}
\langle|\mathcal{M}|^2\rangle=\left(\frac{2mZ\alpha}{q^2}\right)^2
\end{equation*}

The differential cross section is equivalent to the expected probability density.
\begin{equation*}
\frac{d\sigma}{d\Omega}=\langle|\mathcal{M}|^2\rangle=\left(\frac{2mZ\alpha}{q^2}\right)^2
\end{equation*}

Substitute $q^2=2p^2(\cos\theta-1)$ to obtain
\begin{equation*}
\frac{d\sigma}{d\Omega}=\left[\frac{mZ\alpha}{p^2(\cos\theta-1)}\right]^2
\end{equation*}

Noting that
\begin{equation*}
p^2=2mE
\end{equation*}

we can also write
\begin{equation*}
\frac{d\sigma}{d\Omega}=\left[\frac{Z\alpha}{2E(\cos\theta-1)}\right]^2
\end{equation*}

\href{https://georgeweigt.github.io/examples/rutherford-scattering-demo.html}{Eigenmath script}

\end{document}
